\noindent \textit{\textcolor{red!60!black}{
Major Comment 1.\\ 
The authors conduct a placebo test which includes preferred stock and "bond amount outstanding without gold clauses in total liabilities". My impression though is that virtually all corporate bonds of any importance/size at this time included gold clauses. What percent of corporate bonds (by \# and \$) in this sample are 100\% confirmed to not have gold clauses? I would think that number, especially in \$ terms, is tiny.
}}

\bigskip

\noindent Thank you for raising this important point. Your impression is correct: in 1932, about 97\% of the corporate bonds in our sample by count (99\% by outstanding par value) had the gold clause. Put differently, only 3\% of bonds by count and 1\% by dollar value were confirmed to lack gold clauses---numbers that are indeed tiny, as you suspected. In our original submission, we had this information in a footnote, but we agree that such an important piece of information should be featured more prominently in the paper.

To clarify this point, we now mention it on the first page of the introduction: ``When the United States effectively abandoned the gold standard in April 1933, the dollar quickly depreciated against gold, and the existence of these clauses posed a serious threat to most corporate bond issuers: around 97\% of outstanding corporate bonds were subject to the gold clause.'' We hope that explicitly stating the exact number early in the introduction leaves no room for ambiguity and helps frame the economic significance of the gold clause from the outset.

Given this near-universal prevalence of gold clauses in corporate bonds, much of the variation in our measure of gold exposure comes from the composition of firms' liabilities---specifically, the relative share of corporate bonds versus preferred stock. Since preferred shares uniformly lacked gold clauses while nearly all corporate bonds contained them, firms financing with more bonds relative to preferred stock have higher gold clause exposure. As we discuss in our response to Comment 2.d, we have adopted your suggestion to define $d$ as the share of gold clause bonds in long-term debt (bonds plus preferred shares plus bank debt), which directly captures this source of variation. We clarify this when defining our exposure measure in equation (1) and in Section 4.

We have added details on the total number of bonds in our sample and on the number of bonds containing gold clauses. These results are reported in Table 2 and discussed in Section 4 of the revised manuscript. We also emphasize this point when defining gold clause exposure in equation (1), noting that nearly all bonds issued before the leverage uncertainty period included a gold clause.

\bigskip

\noindent \textit{\textcolor{red!60!black}{
On a related note, what percent of all gold clause liabilities, as measured in this paper, come from corporate bonds vs any other liabilities (\# and \$ terms)? Again, my impression from looking through the Moody's Manuals is that virtually all the information on inclusion of gold clauses would be from long-term corporate bonds.
}}

\bigskip

\noindent Your impression is again correct. In our measure $d$, 100\% of gold-denominated liabilities (the numerator) come from long-term corporate bonds having a gold clause. Preferred shares uniformly lacked gold clauses, and while it is technically possible that some bank loans also contained gold clauses, we do not have information on their presence. In any case, bank debt accounts for a negligible fraction of the capital structure of most public firms in our sample. As noted above, this is precisely why variation in $d$ reflects the composition of long-term financing---bonds versus preferred shares---rather than variation in whether bonds contain gold clauses. As we now state in Section 3, ``[w]hile the clause is present in the vast majority of corporate bonds (97\% in our sample), it is entirely absent from preferred shares, another type of long-term obligations commonly used at the time.''


\bigskip

\noindent \textit{\textcolor{red!60!black}{
If I'm mistaken and there is a lot of other variation it would helpful for the authors to clarify that and demonstrate it in the text. It would also suggest that using the \% of long-term corporate bonds with gold clauses, would be a much more natural workhorse empirical design and I would encourage the authors to use that. If, however, as appears to be the case - nearly all the variation in gold clause liabilities is variation in the amount of bonds – that needs to stated and shown explicitly to the reader. To understand the validity of the design it is important to know that the primary source of variation is comparing firms with and without bonds, or with differing amount of bonds relative to their other liabilities, as opposed to firms who happened to include gold clauses or not in their bonds.
}}

\bigskip

\noindent We agree that this point was not sufficiently emphasized in the original submission. As you correctly note, the primary source of variation is comparing firms with differing amounts of bonds relative to their other long-term liabilities---primarily preferred shares---rather than comparing firms whose bonds happened to include or exclude gold clauses. In the revised manuscript, we state this explicitly when defining $d$ in equation (1) and again in Section 4. Specifically, in Section 4 we write:
\begin{quote}
$d_{j,t}$ measures firm $j$'s exposure to the gold clause at time $t$ as the fraction of its long-term liabilities containing the clause:
\begin{equation*}
d_{j,t} = \frac{\text{\normalsize{Gold-denominated long-term liabilities of firm $j$ in year $t$}}}{\text{\normalsize{Long-term liabilities of firm $j$ in year $t$}}}.
\end{equation*}
We define the long-term liabilities of a firm as the sum of the par value of these three components: corporate bonds, preferred shares, and bank loans. To construct $d_{j,t}$, we collect individual security-level information on the presence of the gold clause. As noted in the previous section, the gold clause was included in most corporate bonds (around 97\% of them) and absent from all preferred shares. For bank debt, the manuals do not report exposure to the gold clause. However, at the time of the abrogation of the gold clause bank debt only accounts for 1\% of long-term liabilities [...].
\end{quote}

As discussed in our response to Comment 2.d, we have also adopted your suggestion to use the share of gold clause bonds in long-term debt as our primary specification, which directly reflects this source of variation. We believe these clarifications help readers understand the nature of our identifying variation and the validity of our empirical design.
